\documentclass[11pt]{article}

\usepackage{amsmath}
\usepackage{amssymb}
\usepackage[utf8]{inputenc}
\usepackage[T1]{fontenc}
\usepackage[margin=1in]{geometry}
\usepackage{hyperref}
\hypersetup{
  colorlinks=true,
  linkcolor=blue!50!black,
  urlcolor=blue!50!black
}

% Code style
\usepackage{listings}
\usepackage{xcolor}
\definecolor{codebg}{RGB}{245,245,247}
\lstdefinestyle{py}{
  basicstyle=\ttfamily\small,
  backgroundcolor=\color{codebg},
  frame=single,
  framerule=0pt,
  rulecolor=\color{codebg},
  keywordstyle=\color{blue!60!black}\bfseries,
  commentstyle=\color{green!40!black}\itshape,
  stringstyle=\color{red!40!black},
  showstringspaces=false,
  breaklines=true,
  language=Python,
  xleftmargin=0.5em,
  xrightmargin=0.5em
}
\lstset{style=py}

% Inline code
\newcommand{\code}[1]{\texttt{\small #1}}

\title{HW04 --- Neural Networks\\\large Specification}
\author{Ramaz ML Course 2026}
\date{}

\begin{document}
\maketitle

\noindent
This document is the \textbf{contract} for HW04. It tells you exactly what to
build, what each function and method must do, and the shape of the public API
your code must match. The test suite checks this contract directly --- when
your code matches the spec, the tests pass.

\medskip
\noindent
\textbf{How to use this document:} Read it once end-to-end before opening
any of the Python files. Then implement in roughly the order tests are
written: \code{network.py} first (Dropout, then MLP), then
\code{training.py}, then \code{analysis.py}. Refer back here whenever you
forget a math formula or a return shape.

\medskip
\noindent
All derivations referenced below live in the Lesson 5 LaTeX notes. The spec
points to them rather than re-deriving everything.

\tableofcontents

% ===========================================================================
\section{Architecture overview}
% ===========================================================================

HW04 has three implementation files plus an analysis script:

\begin{center}
\begin{tabular}{ll}
\code{network.py}  & \code{Dropout} (custom nn.Module), \code{MLP} (configurable feedforward net) \\
\code{training.py} & \code{train\_one\_epoch}, \code{validate}, \code{train} (orchestrator) \\
\code{analysis.py} & \code{load\_fashion\_mnist}, \code{build\_model} (you implement);\\
                   & teacher-provided \code{\_\_main\_\_} block runs the ablation
\end{tabular}
\end{center}

\medskip
\noindent
You will implement everything in those three files \emph{above} the
teacher-provided separator in \code{analysis.py}. The training run for Part 2
of the assignment is handled by the \code{if \_\_name\_\_ == "\_\_main\_\_":}
block in \code{analysis.py}, which is released as-is.

\medskip
\noindent
The high-level flow is the canonical PyTorch training loop you saw in lecture:

\begin{center}
\begin{tabular}{ll}
1. & \code{build\_model(config)} returns the model, optimizer, and (optional) scheduler. \\
2. & \code{load\_fashion\_mnist} returns three \code{DataLoader}s. \\
3. & \code{train(...)} loops over epochs, calling \code{train\_one\_epoch} \\
   & and \code{validate}, applying early stopping, restoring the best snapshot. \\
4. & The orchestration block plots curves and prints a results table.
\end{tabular}
\end{center}

% ===========================================================================
\section{\texttt{Dropout} (\texttt{network.py})}
% ===========================================================================

The first thing you implement. \code{Dropout} is a regularization layer that
zeros activations at random during training, scaling the survivors so the
expected activation is unchanged. At evaluation time it does nothing.

\subsection{\texttt{\_\_init\_\_}}

\noindent
\code{def \_\_init\_\_(self, p: float) -> None}

\medskip\noindent
Construct a fresh \code{Dropout} layer with rate $p$. Subclass
\code{torch.nn.Module}; remember to call \code{super().\_\_init\_\_()} before
anything else.

\medskip\noindent
Store:
\begin{itemize}
    \item \code{self.p} --- the dropout probability, as-is.
\end{itemize}

The tests read \code{self.p} after construction, so use this exact attribute
name. \code{p} must lie in $[0, 1)$. If \code{p == 0}, the layer is the
identity in both modes.

\subsection{\texttt{forward}}

\noindent
\code{def forward(self, x: torch.Tensor) -> torch.Tensor}

\medskip\noindent
Apply (or skip) inverted dropout depending on the layer's mode. PyTorch sets
\code{self.training} to \code{True} after \code{model.train()} and
\code{False} after \code{model.eval()} --- read this flag inside
\code{forward} to decide which branch to take.

\medskip\noindent
The \textbf{inverted-dropout} rule (see Lesson 5, Session B, Section
``Dropout'' for the derivation):
\[
    \tilde{x}_i =
    \begin{cases}
        \dfrac{x_i}{1 - p} & \text{with probability } 1 - p \\
        0 & \text{with probability } p
    \end{cases}
\]
when \code{self.training} is \code{True} and $p > 0$; otherwise
$\tilde{x} = x$ (the identity).

\medskip\noindent
The Bernoulli mask is sampled independently per element. Use
\code{torch.rand\_like(x)} as the source of randomness; the autouse seed
fixture in \code{conftest.py} makes this reproducible across test runs.

\begin{itemize}
    \item Input shape: arbitrary (Dropout is shape-preserving).
    \item Output shape: identical to the input.
    \item Expectation: $\mathbb{E}[\tilde{x}] = x$ across the random mask.
\end{itemize}

\subsection*{Invariants the tests will check}

\begin{itemize}
    \item In training mode, with $p = 0.3$ and a large input, the fraction
          of zero entries in the output is approximately $0.3$.
    \item In training mode, every non-zero entry equals $x_i / (1 - p)$.
    \item In eval mode, \code{Dropout(x)} equals \code{x} exactly.
    \item Switching back from \code{eval()} to \code{train()} restores
          dropping behavior.
    \item \code{Dropout} is an \code{nn.Module} subclass.
\end{itemize}

% ===========================================================================
\section{\texttt{MLP} (\texttt{network.py})}
% ===========================================================================

A multi-layer perceptron with configurable depth, hidden widths, and two
optional regularizers. The architecture is exactly the one in Lesson 5
Session A.

\subsection{\texttt{\_\_init\_\_}}

\noindent
\code{def \_\_init\_\_(self, input\_dim: int, hidden\_dims: list[int],}\\
\code{\hspace*{2em}output\_dim: int, dropout\_p: float = 0.0, use\_batchnorm: bool = False) -> None}

\medskip\noindent
Subclass \code{nn.Module} (remember \code{super().\_\_init\_\_()}). Build the
following architecture, in order:

\begin{enumerate}
    \item One \code{nn.Flatten} (so 2-D image inputs like FashionMNIST
          collapse to a single feature dimension).
    \item For each entry $h$ in \code{hidden\_dims}, in order:
    \begin{enumerate}
        \item A \code{nn.Linear} from the previous dimension to $h$.
        \item If \code{use\_batchnorm} is \code{True}: an
              \code{nn.BatchNorm1d(h)}.
        \item An \code{nn.ReLU}.
        \item If \code{dropout\_p > 0}: your custom \code{Dropout(p=dropout\_p)}.
    \end{enumerate}
    \item One final \code{nn.Linear} from the last hidden dimension to
          \code{output\_dim}.
\end{enumerate}

The final layer produces \textbf{raw logits} --- do not apply softmax or any
activation. The loss function (\code{nn.CrossEntropyLoss}) applies the
softmax internally.

\medskip\noindent
Register the layers as submodules so that \code{model.parameters()} returns
all of them. The simplest way is to wrap everything in an
\code{nn.Sequential} and store it as an attribute.

\subsection{\texttt{forward}}

\noindent
\code{def forward(self, x: torch.Tensor) -> torch.Tensor}

\medskip\noindent
Run \code{x} through every layer in order: \code{nn.Flatten}, then the
hidden blocks, then the final linear projection. Return the resulting logit
tensor.

\begin{itemize}
    \item Input shape: \code{(B, 1, 28, 28)} for FashionMNIST, but the
          method must accept any input that \code{nn.Flatten} reduces to
          shape \code{(B, input\_dim)}.
    \item Output shape: \code{(B, output\_dim)}.
\end{itemize}

\subsection*{Invariants the tests will check}

\begin{itemize}
    \item \code{MLP} is an \code{nn.Module} subclass.
    \item Layer counts match the architecture above:
    \begin{itemize}
        \item Exactly one \code{nn.Flatten}.
        \item \code{len(hidden\_dims) + 1} \code{nn.Linear} layers.
        \item \code{len(hidden\_dims)} \code{nn.ReLU} layers.
        \item \code{len(hidden\_dims)} \code{nn.BatchNorm1d} layers when
              \code{use\_batchnorm}, zero otherwise.
        \item \code{len(hidden\_dims)} \code{Dropout} layers when
              \code{dropout\_p > 0}, zero otherwise.
    \end{itemize}
    \item Each \code{nn.Linear}'s \code{in\_features} / \code{out\_features}
          match the adjacent dimensions.
    \item \code{forward} produces \code{(B, output\_dim)} for inputs of
          shape \code{(B, 1, 28, 28)}, for $B \in \{1, 4, 32\}$.
\end{itemize}

% ===========================================================================
\section{\texttt{train\_one\_epoch} (\texttt{training.py})}
% ===========================================================================

\noindent
\code{def train\_one\_epoch(model: nn.Module, loader: DataLoader,}\\
\code{\hspace*{2em}criterion: nn.Module, optimizer: torch.optim.Optimizer,}\\
\code{\hspace*{2em}device: torch.device | str) -> float}

\medskip\noindent
Run one full pass over \code{loader}, updating model parameters. Return the
mean per-example training loss across the epoch.

\medskip\noindent
Algorithm (one iteration of the standard PyTorch training loop):

\begin{enumerate}
    \item Switch the model to training mode (so \code{Dropout} and
          \code{BatchNorm1d} behave correctly).
    \item For each \code{(x, y)} batch in \code{loader}:
    \begin{enumerate}
        \item Move the batch tensors onto \code{device}.
        \item Zero the optimizer's accumulated gradients.
        \item Feed the batch through the model to obtain its logit predictions.
        \item Compute the scalar loss with \code{criterion}.
        \item Run the backward pass on the loss.
        \item Step the optimizer.
        \item Accumulate the loss (weighted by batch size) for the
              epoch-mean.
    \end{enumerate}
    \item Return the loss totalled across all examples, divided by the
          number of examples.
\end{enumerate}

\begin{itemize}
    \item Return value is a Python \code{float}, not a tensor.
    \item Order of operations: zero-grad before forward, step after
          backward. Reversing these silently breaks training.
\end{itemize}

% ===========================================================================
\section{\texttt{validate} (\texttt{training.py})}
% ===========================================================================

\noindent
\code{def validate(model: nn.Module, loader: DataLoader, criterion: nn.Module,}\\
\code{\hspace*{2em}device: torch.device | str) -> tuple[float, float]}

\medskip\noindent
Evaluate the model on \code{loader} without updating parameters. Return
\code{(loss, accuracy)} as Python floats.

\medskip\noindent
Algorithm:

\begin{enumerate}
    \item Switch the model to eval mode (so \code{Dropout} is the identity
          and \code{BatchNorm1d} uses its running statistics).
    \item Open a \code{torch.no\_grad()} context (prevents PyTorch from
          building the computation graph; saves memory and time).
    \item For each \code{(x, y)} batch in \code{loader}:
    \begin{enumerate}
        \item Move the batch tensors onto \code{device}.
        \item Feed the batch through the model to obtain logits.
        \item Compute the loss with \code{criterion}.
        \item Take the argmax of the logits along the class dimension to
              get a predicted class index per example.
        \item Accumulate the loss (weighted by batch size) and the count
              of correct predictions.
    \end{enumerate}
    \item Return the mean loss and the mean accuracy across all examples
          as a two-tuple.
\end{enumerate}

\begin{itemize}
    \item Both return values are Python \code{float}s.
    \item Accuracy is always in $[0, 1]$.
    \item The function must \textbf{not} modify the model parameters.
\end{itemize}

% ===========================================================================
\section{\texttt{train} (\texttt{training.py})}
% ===========================================================================

\noindent
\code{def train(model: nn.Module, train\_loader: DataLoader, val\_loader: DataLoader,}\\
\code{\hspace*{2em}optimizer: torch.optim.Optimizer, criterion: nn.Module,}\\
\code{\hspace*{2em}max\_epochs: int, patience: int, device: torch.device | str,}\\
\code{\hspace*{2em}scheduler: torch.optim.lr\_scheduler.LRScheduler | None = None) -> dict}

\medskip\noindent
The full training orchestrator. Runs the train/validate cycle, drives the
optional LR scheduler, applies inline early stopping, keeps the
best-validation parameters in memory, and restores them before returning.

\medskip\noindent
Algorithm:

\begin{enumerate}
    \item Initialize:
    \begin{itemize}
        \item An empty history dict with keys \code{train\_loss},
              \code{val\_loss}, \code{val\_acc}, \code{best\_epoch}. The
              first three start as empty lists; \code{best\_epoch} starts
              at $0$.
        \item A best-loss-so-far variable, set to $+\infty$.
        \item A deep copy of the current \code{model.state\_dict()} as the
              best snapshot.
        \item An ``epochs since improvement'' counter, set to $0$.
    \end{itemize}
    \item For epoch $= 0$ up to \code{max\_epochs - 1}:
    \begin{enumerate}
        \item Call \code{train\_one\_epoch}; append the returned loss to
              \code{history['train\_loss']}.
        \item Call \code{validate}; append the loss and accuracy to
              \code{history['val\_loss']} and \code{history['val\_acc']}.
        \item If \code{scheduler is not None}, call \code{scheduler.step()}.
        \item If the new val loss is strictly less than the best-so-far:
              update best, deep-copy \code{model.state\_dict()} as the new
              best snapshot, set \code{history['best\_epoch']} to the
              current epoch, reset the counter to $0$.
        \item Otherwise: increment the counter. If it reaches
              \code{patience}, break out of the loop (early stop).
    \end{enumerate}
    \item Restore the best snapshot via \code{model.load\_state\_dict(...)}.
    \item Return the history dict.
\end{enumerate}

\medskip\noindent
\textbf{Deep copy matters.} \code{model.state\_dict()} returns a \emph{view}
of the current parameters; if you assign it directly to your ``best''
variable, that view will be updated on every subsequent gradient step
(meaning ``best'' tracks the latest, not the best). Use \code{copy.deepcopy}
on the state dict.

\begin{itemize}
    \item Return type: \code{dict} with the four keys above. \code{train\_loss},
          \code{val\_loss}, \code{val\_acc} are \code{list[float]} of equal
          length. \code{best\_epoch} is an \code{int}.
    \item If early stopping does not trigger, the lists have length
          \code{max\_epochs}.
    \item After \code{train} returns, the model's parameters reflect the
          best-validation epoch, not the last epoch.
\end{itemize}

% ===========================================================================
\section{Analysis helpers (\texttt{analysis.py})}
% ===========================================================================

\code{analysis.py} contains two helper functions that the teacher-provided
\code{\_\_main\_\_} block uses to run the ablation study. The \code{\_\_main\_\_}
block itself is released as-is below the marked separator --- do not modify
it. Match the helper signatures exactly.

\subsection{\texttt{load\_fashion\_mnist}}

\noindent
\code{def load\_fashion\_mnist(batch\_size: int = 128)}\\
\code{\hspace*{2em}-> tuple[DataLoader, DataLoader, DataLoader]}

\medskip\noindent
Load FashionMNIST and return train, val, and test \code{DataLoader}s as a
three-tuple. This is a utility helper, so the spec names the relevant
library entry points directly.

\medskip\noindent
Steps:
\begin{enumerate}
    \item Build a transform pipeline of \code{transforms.ToTensor()} followed
          by \code{transforms.Normalize((0.2860,), (0.3530,))} (the
          FashionMNIST channel mean and std).
    \item Load \code{torchvision.datasets.FashionMNIST(root="./data",
          train=True, download=True, transform=transform)}.
    \item Compute split sizes: $80\%$ for training, $10\%$ for validation,
          $10\%$ for test, using integer multiplication on the total length.
          The remainder (if any) goes to the test split.
    \item Build a \code{torch.Generator().manual\_seed(42)} so the split is
          reproducible across runs.
    \item Call \code{torch.utils.data.random\_split} with the three sizes
          and the generator to split the dataset.
    \item Wrap each split in a \code{DataLoader} with the given
          \code{batch\_size}. Shuffle the train loader; leave val and test
          unshuffled. Use \code{num\_workers=0}.
\end{enumerate}

\noindent
Return order: \code{(train\_loader, val\_loader, test\_loader)}. The split
sizes for the FashionMNIST training set of 60,000 examples are
$48000 / 6000 / 6000$.

\subsection{\texttt{build\_model}}

\noindent
\code{def build\_model(config: str)}\\
\code{\hspace*{2em}-> tuple[nn.Module, torch.optim.Optimizer,}\\
\code{\hspace*{2em}~~~~torch.optim.lr\_scheduler.LRScheduler | None]}

\medskip\noindent
Assemble the model, optimizer, and (optional) scheduler trio for one of the
four ablation configurations. The shared architecture is:
\code{MLP(input\_dim=784, hidden\_dims=[256, 128], output\_dim=10, ...)};
constants \code{INPUT\_DIM}, \code{HIDDEN\_DIMS}, \code{OUTPUT\_DIM} are
defined at the top of \code{analysis.py}.

\medskip\noindent
The four valid \code{config} values and what they each build:

\begin{center}
\begin{tabular}{lcccc}
\code{config}  & \code{dropout\_p} & \code{use\_batchnorm} & \code{weight\_decay} & scheduler \\
\hline
\code{"baseline"}      & 0.0 & False & 0       & None \\
\code{"dropout"}       & 0.3 & False & 0       & None \\
\code{"batchnorm"}     & 0.0 & True  & 0       & None \\
\code{"both\_schedule"} & 0.3 & True  & $10^{-4}$ & StepLR \\
\end{tabular}
\end{center}

\medskip\noindent
For all four configs, use \code{torch.optim.SGD} with \code{lr=0.1}. When a
scheduler is requested (\code{"both\_schedule"} only), use
\code{torch.optim.lr\_scheduler.StepLR(optimizer, step\_size=10, gamma=0.5)}.
For an unknown \code{config} string, raise a \code{ValueError}.

\begin{itemize}
    \item Return order: \code{(model, optimizer, scheduler)}, with
          \code{scheduler} being \code{None} when not applicable.
    \item All three configurations without a scheduler must return
          \code{None} for the third element --- not a no-op object.
\end{itemize}

% ===========================================================================
\section{How to read the test file}
% ===========================================================================

\code{test\_network.py} is organized in three layers --- work on them in order:

\subsection*{\texttt{@pytest.mark.unit} --- one piece at a time (35 pts)}

Per-method correctness for \code{Dropout} and \code{MLP}. Get these green
first; they verify the building blocks of \code{network.py} independently
of any training.

\subsection*{\texttt{@pytest.mark.integration} --- training pipelines (30 pts)}

End-to-end behavior for \code{train\_one\_epoch}, \code{validate}, and
\code{train}. These tests use tiny synthetic datasets so they run in seconds.

\subsection*{\texttt{@pytest.mark.analysis} --- analysis helpers (35 pts)}

Tests for \code{load\_fashion\_mnist} (correct splits, shapes, dtypes) and
\code{build\_model} (each of the four configs returns the right architecture
and optimizer settings). The first run will download FashionMNIST
(~30 MB); subsequent runs use the cache in \code{./data/}.

\medskip
\noindent
Run a single layer with:

\begin{lstlisting}[language=bash]
uv run pytest -m unit         # just unit tests
uv run pytest -m integration  # just integration tests
uv run pytest -m analysis     # just analysis tests
\end{lstlisting}

\medskip
\noindent
Each test class is \textbf{all-or-nothing} --- every test in a class must
pass to earn the points for that class.

% ===========================================================================
\section{Restrictions}
% ===========================================================================

\begin{itemize}
    \item \textbf{No \code{torch.nn.Dropout}} anywhere in \code{network.py}.
          Implementing dropout from scratch is the substantive learning goal
          of §2. The \code{conftest.py} parses your source with \code{ast}
          and fails the run if it finds any reference to \code{nn.Dropout},
          \code{torch.nn.Dropout}, or \code{from torch.nn import Dropout}.
    \item \textbf{No \code{sklearn.neural\_network}.} Build the MLP from
          \code{torch.nn} primitives.
\end{itemize}

You \emph{may} use freely:
\begin{itemize}
    \item \code{torch.autograd}, \code{loss.backward()}, \code{.requires\_grad}
          --- this is the first assignment where these are allowed.
    \item \code{nn.Linear}, \code{nn.ReLU}, \code{nn.BatchNorm1d},
          \code{nn.Flatten}, \code{nn.Sequential}.
    \item \code{torch.optim.SGD}, \code{torch.optim.lr\_scheduler.StepLR}.
    \item \code{torch.utils.data.DataLoader}, \code{torch.utils.data.random\_split}.
    \item \code{torchvision.datasets.FashionMNIST}, \code{torchvision.transforms}.
\end{itemize}

\end{document}
